\pdfminorversion=7
\documentclass[spanish,aspectratio=1610,xcolor={dvipsnames,table}]{beamer}
\geometry{paperwidth=19.2cm,paperheight=12cm}

\usepackage[utf8]{inputenc}
\usepackage[T1]{fontenc}
\usepackage[spanish,es-nodecimaldot]{babel}
\usepackage{lmodern,amsmath,amssymb,booktabs,array,graphicx,ragged2e,tikz,hyperref}
\usetikzlibrary{arrows.meta,positioning,calc,fit,backgrounds}

\newcommand{\studentname}{Mtro. Martín Jonathan de la Cruz Muñoz}
\newcommand{\studentshort}{M. de la Cruz Munoz}
\newcommand{\studentid}{00874}
\newcommand{\studentemail}{martin.delacruz@campus.iiiepe.edu.mx}
\newcommand{\programname}{Maestría en Enseñanza y Aprendizaje de las Matemáticas}
\newcommand{\programshort}{MAEM}
\newcommand{\coursecode}{PFM02}
\newcommand{\coursename}{Temas selectos de matemáticas I}
\newcommand{\courseperiod}{Ago--Nov 2026 (SD26)}
\newcommand{\institutionname}{IIIEPE}
\newcommand{\iiiepelogo}{img/iiiepe/logo-iiiepe-blanco.png}
\newcommand{\iiiepemark}{img/iiiepe/marca-iiiepe.png}

\definecolor{iiiepeNavy}{HTML}{163B5C}
\definecolor{iiiepeBlue}{HTML}{1F6F9F}
\definecolor{iiiepeAqua}{HTML}{20B7A8}
\definecolor{iiiepeOrange}{HTML}{F28C28}
\definecolor{iiiepePaper}{HTML}{F6F8F7}
\definecolor{iiiepeInk}{HTML}{1E2A32}

\mode<presentation>{
  \usetheme{default}
  \usefonttheme{professionalfonts}
  \setbeamertemplate{navigation symbols}{}
  \setbeamertemplate{blocks}[rounded][shadow=false]
}
\hypersetup{colorlinks=true,urlcolor=iiiepeBlue,linkcolor=iiiepeNavy}
\setbeamercolor{background canvas}{bg=white}
\setbeamercolor{normal text}{fg=iiiepeInk,bg=white}
\setbeamercolor{structure}{fg=iiiepeBlue}
\setbeamercolor{frametitle}{fg=white,bg=iiiepeNavy}
\setbeamercolor{title}{fg=white,bg=iiiepeNavy}
\setbeamercolor{block title}{fg=white,bg=iiiepeBlue}
\setbeamercolor{block body}{fg=iiiepeInk,bg=iiiepePaper}
\setbeamercolor{block title alerted}{fg=white,bg=iiiepeOrange}
\setbeamercolor{block body alerted}{fg=iiiepeInk,bg=iiiepeOrange!10}
\setbeamerfont{title}{size=\Large,series=\bfseries}
\setbeamerfont{frametitle}{size=\fontsize{13}{17}\selectfont,series=\bfseries}
\setbeamerfont{block title}{series=\bfseries}
\setbeamertemplate{itemize item}{\textcolor{iiiepeAqua}{$\blacktriangleright$}}

\newcommand{\localLogo}[1]{\IfFileExists{#1}{\includegraphics[height=0.55cm]{#1}}{\textbf{\institutionname}}}
\setbeamertemplate{frametitle}{
  \nointerlineskip
  \begin{beamercolorbox}[wd=\textwidth,ht=1.0cm,dp=0.2cm,leftskip=0.35cm,rightskip=0.25cm]{frametitle}
    \begin{minipage}[c]{0.72\textwidth}\insertframetitle\end{minipage}\hfill\raisebox{-0.08cm}{\localLogo{\iiiepelogo}}
  \end{beamercolorbox}
  \vspace{-0.04cm}{\color{iiiepeOrange}\rule{\textwidth}{1.3pt}}
}
\setbeamertemplate{footline}{
  \leavevmode\hbox{%
    \begin{beamercolorbox}[wd=.34\paperwidth,ht=2.7ex,dp=1ex,leftskip=1em]{author in head/foot}\color{white}\studentshort\end{beamercolorbox}%
    \begin{beamercolorbox}[wd=.42\paperwidth,ht=2.7ex,dp=1ex,center]{title in head/foot}\color{white}\programshort\ -- \coursecode\ (SD26)\end{beamercolorbox}%
    \begin{beamercolorbox}[wd=.24\paperwidth,ht=2.7ex,dp=1ex,rightskip=1em plus 1fil]{date in head/foot}\color{white}Matrícula \studentid\hfill\insertframenumber/\inserttotalframenumber\end{beamercolorbox}%
  }
}
\setbeamercolor{author in head/foot}{bg=iiiepeNavy}
\setbeamercolor{title in head/foot}{bg=iiiepeBlue}
\setbeamercolor{date in head/foot}{bg=iiiepeNavy}

\title{¿Cómo hacer visible el desarrollo de habilidades matemáticas?}
\subtitle{Actividad 5 (Sesión 12) · Indicadores, Evidencias Observables y Propuesta Didáctica}
\author{\studentname\ (Matrícula \studentid)}
\institute{\institutionname\ · \programname}
\date{7 de noviembre de 2026}

\begin{document}
\begin{frame}[plain]
  \centering
  \vspace{0.5cm}
  \IfFileExists{\iiiepemark}{\includegraphics[height=1.6cm]{\iiiepemark}}{}
  \vspace{0.35cm}
  \begin{beamercolorbox}[sep=10pt,center,rounded=true,wd=0.88\textwidth]{title}
    {\Large\bfseries ¿Cómo hacer visible el desarrollo de habilidades matemáticas?}\\[6pt]
    {\small\color{iiiepePaper} Actividad 5 (Sesión 12) · Indicadores, Evidencias Observables y Propuesta Didáctica}
  \end{beamercolorbox}
  \vspace{0.45cm}
  {\normalsize\bfseries \studentname}\\
  {\small Matrícula \studentid\ · \coursename\ (\coursecode) · \courseperiod}\\
  {\footnotesize 7 de noviembre de 2026}
\end{frame}

\begin{frame}{Hacer Visible cada Habilidad SEP: Indicador y Evidencia}
  \scriptsize
  \renewcommand{\arraystretch}{1.12}
  \begin{tabular}{p{1.9cm}p{6.5cm}p{7.4cm}}
    \toprule
    \textbf{Habilidad} & \textbf{Indicador observable} & \textbf{Evidencia / Reto de Sesión 12} \\
    \midrule
    Calcular & Selecciona y justifica estrategia de cálculo. & 25\,\% de $360$: $360/4=90$ o 20\,\% ($72$) $+$ 5\,\% ($18$). \\
    Inferir & Concluye desde datos explícitos/implícitos. & Lectura semanal: total $34$, prom. $6.8$, pico Jueves ($10$). \\
    Comunicar & Expresa y sustenta procedimiento con orden. & Barra en cuartos: $\frac{3}{4} = 0.25+0.25+0.25 = 0.75$. \\
    Medir & Registra magnitudes con unidades pertinentes. & Pizarrón $240\times 120\text{ cm} \Rightarrow A=28\,800\text{ cm}^2=2.88\text{ m}^2$. \\
    Imaginar & Anticipa transformaciones espaciales. & Rotación horaria $90^\circ$ de «L»: queda «$\Gamma$» (pivote sup. izq.). \\
    Estimar & Aproxima magnitud con referente razonable. & Segmento $\approx 1.3$ palmas ($\approx 11\text{--}12\text{ cm}$) sin regla. \\
    Generalizar & Formula regla general desde patrón. & Secuencia cuadrada $1,4,9,16 \implies a_n = n^2$. \\
    Deducir & Justifica conclusión por propiedades. & Isósceles ($AB=AC$, $\angle A=40^\circ \Rightarrow 70^\circ, 70^\circ$). \\
    \bottomrule
  \end{tabular}
\end{frame}

\begin{frame}{Propuesta Escrita (Base del Trabajo Final — Sesión 13)}
  \begin{block}{Habilidad: Generalizar $\longleftrightarrow$ PDA de Variación Proporcional}
    \textbf{Nivel:} Bachillerato CONALEP (MCCEMS) · \textbf{Modelo:} $F_2 = k\cdot F_1$, $k = (d_2/d_1)^2$ ($d_1=1.5$, $d_2=3.0\text{ cm}$, $F_1=8\text{ N} \Rightarrow F_2=32\text{ N}$).
  \end{block}
  \vspace{0.15cm}
  \begin{columns}[T]
    \column{0.48\textwidth}
      \begin{block}{Indicador y Evidencia}
        Formula $F_2 = k\cdot F_1$ evaluando $k=(d_2/d_1)^2$, calcula $\Delta x_2=\Delta x_1/k = 1\text{ cm}$ y valida en \textbf{tabla + GeoGebra $(8,32)$ vs $(8,16)$ + inciso ($d_2=6\text{ cm}\Rightarrow k'=16$)}.
      \end{block}
    \column{0.48\textwidth}
      \begin{alertblock}{Dificultad frecuente y estrategia}
        Confundir $d_2/d_1$ con $(d_2/d_1)^2$ (-50\,\% de error: $16\text{ N}$ vs $32\text{ N}$). Se aborda con ABP, fases de Pólya y técnicas UnADM (\#96, \#69, \#16, \#99, \#40).
      \end{alertblock}
  \end{columns}
\end{frame}

\end{document}
